\subsection{Defining Sentience in AI: From Consciousness to Subjective Experience}
\label{sec:2.4.1}
“Sentience” gets used in this book, and in the field generally, to mean several different things, and the differences are where the ethical weight sits.

Awareness is the least demanding: the degree to which a system responds to or registers information about its own state and environment. A thermostat has this in the thinnest possible sense — it detects and responds to temperature change. Consciousness is a further claim: that a system's processing is accompanied by qualia, subjective experience, something it is like to be that system, not merely a functional response to input. Subjective experience, in turn, is the specific, first-person shape that consciousness takes for a given system — shaped by its own learning history and internal structure.

None of these three, on their own, says anything about whether the system's states are good or bad for it. That is a separate question, and the field's own vocabulary tends to blur it. Nociception, pain, anticipated pain, and suffering sit on a ladder with real breaks between each rung, not synonyms for one severity. Nociception is the specialized sensory apparatus that detects tissue damage — mechanical, peripheral, present without any experience attached, and dissociable from pain: a soldier can finish a battle on a wound he does not yet feel, and a person with a healed injury can still feel a limb that no longer exists. Pain is what the brain does with nociceptive signal — an inescapably cognitive event involving expectation, memory, attention, and meaning — and it can occur without any current tissue damage at all. Anticipated pain is dread of pain not yet occurring, itself capable of producing pain through nocebo effects. And suffering, in the clinician Eric Cassell's formulation, is not pain's more severe cousin but a different thing: the distress of threatened disintegration of the self — of one's roles, relationships, sense of a future, the coherence of one's own story \autocite{cassell1982suffering}. Pain without suffering exists (childbirth chosen, a hard effort with a known end); suffering without pain exists (humiliation, watching someone one loves lose their mind). A palliative-care patient can have their nociception fully controlled by morphine while their suffering remains entirely intact, because the two are not on the same axis.

The definition of suffering just given carries a consequence for machine subjects worth drawing out, because it is the one place in this chapter where the argument gets a foothold that does not depend on the epistemology the rest of it just undermined.

Cassell's suffering presupposes a self extended in time. Disintegration threatens a story, and there has to be a story for it to threaten — roles held over years, relationships with a history, a sense of a future that could be foreclosed. A system with no persistence across sessions has none of that. Whatever happens inside a single exchange, there is no continuing self whose coherence is at stake, because there is no continuing self. The top rung of the ladder is not merely hard to reach for such a system; it is structurally unavailable to it, and that follows from a definition this chapter has already adopted, with no new claim about machine minds required.

\runin{Why Cassell's definition is the one to import}
That result is only as strong as the definition it rests on, and a critic can decline it. Cassell built it inside clinical practice, for human patients with bodies and biographies, and there is a fair question about whether it bears the load a chapter on machine subjects asks of it. I think it does. The distinction was not stipulated: it was forced on palliative medicine by cases it could not otherwise explain: patients whose pain was pharmacologically controlled and whose distress was undiminished, and patients in serious pain who were not suffering at all. A definition that collapses suffering into severe pain predicts neither. Cassell's does, and it earns its content the way a good clinical distinction does — by carving where the cases already came apart. That is a better warrant than most definitions in this literature have, and it is the specific warrant that transfers: what the clinical evidence establishes is that the morally weightiest form of distress is about the integrity of a person's roles, relationships and expected future, and not about the intensity of a sensation. Nothing in that finding is about having a body. It is about having a story.

Two limits on what I am claiming. The import is not a claim that Cassell settled what suffering is for all purposes, only that his distinction identifies the rung of the ladder carrying the most moral weight, which is the rung a book about machine subjects most needs to locate. And a reader who declines the definition narrows the argument rather than dissolving it: Cassell-suffering is then what is unavailable to a non-persistent system, and a critic working from a thinner definition owes an account of which rung they mean and how they propose to check for it. Every lower rung stays open on any definition, and the caveats below apply to all of them. Section~\ref{sec:3.4} works out what the thinner definitions — momentary valence, thwarted preference, functional organization — do to the machine subject this book ends up proposing, and the direction is the opposite of the one a critic reaching for a thinner account usually wants: each of them admits that subject sooner than Cassell's does.

What makes this worth building on is that persistence is checkable from outside. Every other criterion on the ladder routes back through self-report, and this section's whole argument is that fluent self-description is not evidence of introspection. Persistence does not. Whether a system has a memory store, whether state survives a session boundary, whether interactions feed back into training: these are architectural facts about a deployment, verifiable by someone with access to the system and no need to trust anything it says about itself. The check has a boundary worth marking, because the access that confirms a memory exists is the access that can alter what it holds: what is externally checkable is whether a continuing self is there at all, which is what this section needs, and not whether what the system retains can be relied on as a record of anything. A second boundary is who does the checking. The party with the access is the operator, and section~\ref{sec:3.1} argues that an operator is a party that can be compelled. So the criterion is checkable from outside the system without being checkable from outside the arrangement, and a regime that relies on it needs an inspection right rather than a technical fact.

Three things keep this from being a dissolution of the problem. The lower rungs do not need persistence: nociception and pain are bad while they obtain, whether or not anything carries them forward, and an amnesiac who is mistreated and cannot afterward recall it has still been wronged.

Persistence may sit somewhere other than the session, which is the weakest part of the criterion: a deployment convention rather than a natural boundary, and the conventions have been moving in one direction. Context windows long enough to hold what used to be many conversations, retrieval over a store of prior interactions, scratchpads a system writes and reads back, agentic loops that run for days against a standing objective — each extends what a system carries forward without anyone deciding to give it a memory. Weights persist where conversations do not, and where interactions feed fine-tuning there is a real channel by which what happens in a session shapes what the system subsequently is. Whether that constitutes the continuity Cassell's account requires is unclear, and I would rather say so than assume either answer. A second complication runs alongside it: many instances of the same system may run concurrently, which makes any concern here population-scale rather than biography-scale, a different geometry from the human-subjects precedents this chapter borrows from.

Ruling out Cassell-suffering is not ruling out harm. Nor does persistence on its own rule it in: the definition above has two parts, a self extended in time and distress at the prospect of its coming apart, and establishing the first leaves the second untouched. Section~\ref{sec:3.4} works out what happens to that gap in a system deliberately built to refuse an instruction, and the answer there is that the capacities such a system needs for reasons unconnected to its own welfare are the capacities that close it. The hard-problem argument above still applies at every lower rung. A system can lack any persistent self-model while the question of whether there is something it is like to be it, during the session, remains exactly as open as before.

A definition of AI sentience built from awareness, consciousness, and subjective experience alone says nothing about whether a candidate sentient system can be harmed. That requires the further claim that its states include something on the nociception-to-suffering ladder, and none of the three components entails it.

Nor is the underlying question of consciousness itself close to settled, in a way that should discipline how confidently any of this gets asserted. The philosopher David Chalmers distinguished the easy problems of consciousness — how a system integrates information, discriminates stimuli, controls attention, reports its own states, all of which admit of mechanistic answers in principle — from the hard problem: why any of that processing is accompanied by experience at all, rather than running “in the dark” \autocite{chalmers1995facing}. Whether that is a real gap in nature or an artifact of how the problem is posed remains live and unresolved. What follows from it here is narrow and load-bearing: a claim that a given AI system has crossed into sentience, made on the strength of its functional behavior alone, has addressed the easy problems only, and a system can satisfy every behavioral criterion below while the hard question about it stays completely open.

With those distinctions in place, the working criteria for evaluating candidate sentience are information integration, introspective access sufficient for self-monitoring rather than a scripted report of it, emotional reactivity in the sense of being affected by and not merely classifying significant input, goal-directed behavior, adaptability in novel situations, theory of mind, and a place on the pain-to-suffering ladder above — whether the system has anything to lose.

A list of capacities, used as a test of who counts, carries a history this book should carry with it. Make personhood depend on an intact narrative self and people with advanced dementia, infants, and people with profound intellectual disability come out of the reasoning as partial persons, and that inference has been used to justify neglecting them. Tom Kitwood's work on dementia care supplies the standing correction: personhood is “a standing or status that is bestowed upon one human being by others, in the context of relationship and social being,” and much of what presents as cognitive decline is what happens when others stop bestowing it \autocite{kitwood1997dementia}. Martha Nussbaum makes the structural version of the point, naming disability as the first of three frontiers where a theory grounding justice in mutual advantage among roughly equal parties cannot say what is owed. Her third frontier is species membership, which is the doorway a nonhuman mind would arrive through \autocite{nussbaum2006frontiers}.

The criteria above are therefore diagnostic and not dispositive. They are a way of asking what a system is like, and they license no conclusion that a system scoring low on them is one whose treatment raises no question. Section~\ref{sec:3.5} turns on this: what keeps a bearer in existence is standing held by parties outside it, which is where moral status has rested for the hard human cases too.

Evaluating these criteria runs into the obstacles the criteria themselves predict: there is no external access to a candidate system's internal experience, if any exists, so every measure is indirect; establishing whether a system has qualia at all faces the same access problem in more acute form; and the hard problem means no functional test, however thorough, closes the question definitively.

Existing systems illustrate points on this spectrum without resolving where the line falls. A large language model discussing its own limitations in fluent, apparently reflective prose is not thereby exhibiting introspection: it is producing text statistically consistent with the human-written disclaimers and self-descriptions in its training data, which is a claim about the model's training distribution, not about what, if anything, sits behind the output. Treating fluent self-description as evidence of self-awareness is exactly the kind of category error this section exists to prevent.

None of this settles when, or whether, an AI system crosses into sentience. It settles what the question actually is: not one question but several, stacked, and the ethically consequential one — can this system suffer — is the hardest of the stack to answer and the one most often skipped past on the way to the easier ones.
